\section{Expected Dynamic Behavior}

\subsection*{Qualitative shape of $P(t)$}

Even without the innate brake, the 4-state chain produces a characteristic delayed rise $\to$ peak $\to$ decay, because antigen production requires sequential progression through multiple steps (depot $\to$ uptake $\to$ cytosol $\to$ protein). This is consistent with many reported mRNA--LNP expression profiles where maximal signal occurs after several hours and then declines. \parencite{pardi2015expressionKinetics}

The model also naturally produces:
\begin{itemize}
    \item \textbf{Delay before antigen accumulation}: controlled mainly by $1/(k_{\mathrm{up}}+k_{\mathrm{clr},L})$ and $1/(k_{\mathrm{rel}}+k_{\mathrm{loss},E})$.
    \item \textbf{Peak timing}: controlled by the slowest upstream rates plus protein turnover.
    \item \textbf{Persistence/tail}: controlled by the slower of (i) depot persistence, (ii) cytosolic mRNA decay, and (iii) protein clearance.
\end{itemize}

Adding the optional $I(t)$ module can yield sharp drops after an early peak, consistent with cases where strong innate responses produce translational inhibition. \parencite{Huysmans2019saRNALNPSkin}

\subsection*{Analytical insights}

Because the main model is linear, several ``global'' outcomes can be derived without simulation.

\subsubsection*{Escape and uptake fractions}

Define two dimensionless efficiencies:
\begin{align}
 f_{\mathrm{up}} &= \frac{k_{\mathrm{up}}}{k_{\mathrm{up}} + k_{\mathrm{clr},L}} \, \qquad \text{(fraction of depot that becomes internalized)}, \\
 f_{\mathrm{esc}} &= \frac{k_{\mathrm{rel}}}{k_{\mathrm{rel}} + k_{\mathrm{loss},E}} \, \qquad \text{(fraction of internalized mRNA that becomes cytosolic)}.
\end{align}

Reported cytosolic delivery efficiencies for LNP systems motivate $f_{\mathrm{esc}} \ll 1$ in many regimes (order $0.01$--$0.03$). \parencite{Pei2025EndosomalEscapeLNP}

\subsubsection*{Total cytosolic mRNA exposure}

Integrate $\mathrm{d}M/\mathrm{d}t = k_{\mathrm{rel}}E - k_{\mathrm{deg}}M$ from $0$ to $\infty$ (with $M(0)=M(\infty)=0$):
\begin{equation}
\int_0^{\infty} k_{\mathrm{rel}}E(t)\,\mathrm{d}t = k_{\mathrm{deg}}\int_0^{\infty} M(t)\,\mathrm{d}t.
\end{equation}

Similarly, conservation through the upstream chain gives total cytosolic ``arrivals'':
\begin{equation}
\int_0^{\infty} k_{\mathrm{rel}}E(t)\,\mathrm{d}t = D_0 f_{\mathrm{up}} f_{\mathrm{esc}}.
\end{equation}

Therefore,
\begin{equation}
\int_0^{\infty} M(t)\,\mathrm{d}t = \frac{D_0 f_{\mathrm{up}} f_{\mathrm{esc}}}{k_{\mathrm{deg}}}.
\end{equation}

Interpretation: cytosolic mRNA exposure scales multiplicatively with (i) how much dose is internalized, (ii) how much escapes, and (iii) the cytosolic lifetime.

\subsubsection*{Antigen exposure (AUC) as a product of mechanisms}

Integrate $\mathrm{d}P/\mathrm{d}t = k_{\mathrm{tl}}M - k_{\mathrm{clr}}P$ (with $P(0)=P(\infty)=0$):
\begin{equation}
 k_{\mathrm{tl}}\int_0^{\infty} M(t)\,\mathrm{d}t = k_{\mathrm{clr}}\int_0^{\infty} P(t)\,\mathrm{d}t.
\end{equation}

So:
\begin{equation}
\mathrm{AUC}_P \equiv \int_0^{\infty} P(t)\,\mathrm{d}t = D_0 f_{\mathrm{up}} f_{\mathrm{esc}}\frac{k_{\mathrm{tl}}}{k_{\mathrm{deg}}k_{\mathrm{clr}}}.
\end{equation}

Interpretation:
\begin{itemize}
    \item If endosomal escape is 1--3\%, it can dominate the overall scaling even if translation is strong. \parencite{Pei2025EndosomalEscapeLNP}
    \item Improvements in mRNA chemistry/purity that reduce innate activation can increase the effective $k_{\mathrm{tl}}$ by preventing translation shutdown, linking immunogenicity to expression kinetics. \parencite{Nelson2020mRNAChemistryInnateImmunity}
\end{itemize}

\subsection*{Peak vs persistence trade-off}

In this linear chain, higher persistence can be achieved by reducing $k_{\mathrm{clr},L}$, $k_{\mathrm{deg}}$, or $k_{\mathrm{clr}}$, but those changes often shift the peak later and broaden the curve (longer transit/accumulation). Conversely, aggressive early uptake/release (large $k_{\mathrm{up}}$, effectively larger early flux) can create early high peaks but short duration if clearance/decay rates are also high.
